\documentclass[11pt]{article}
\usepackage[margin=1in]{geometry}
\usepackage{hyperref}
\title{Explainable AI with Permutation Importance}
\author{Devis W. Saputra}
\begin{document}
\maketitle

\begin{abstract}
I train a 450-tree Random Forest on the Wisconsin Diagnostic Breast Cancer dataset and use repeated permutation importance on held-out data to inspect which features the fitted model depends on.
\end{abstract}

\section{Question}
Which input features most affect the held-out performance of a strong nonlinear classifier?

\section{Data and Method}
I use a stratified 75/25 split. The Random Forest has 450 trees and random state 42. Permutation importance is calculated on the test set with 16 repeats per feature.

\section{Results}
The held-out ROC-AUC is 0.9945. The largest measured permutation effects come from worst texture and worst smoothness, followed by several concavity-related measurements.

\section{Interpretation}
Permutation importance describes model dependence, not biological causality. Correlated features can also share information and reduce one another's measured importance.

\section{Limitations}
The ranking comes from one fitted model and one split. A stronger study would repeat training, inspect explanation stability, and compare with other explanation methods.

\section{Reproduction}
Install \texttt{requirements.txt} and run \texttt{python src/run\_experiment.py}.

\bibliographystyle{plain}
\bibliography{references}
\end{document}
